\scanpage{64}
% loc: 1–4, top
On trouve une représentation fidèle de $\mathrm{SL}(2,\mathbb Z)/{\pm1}$ \deleted{dans} (engendré par $\sigma$ et $\rho_f$ mod Centre de l'Alg. d'Azumaya.) \deleted{\unsure{du gr. dans la place|Supp. dans le plus}}. Considérons

% loc: 5–8
\deleted{Soit} la représentation de la carte triangulaire universelle, dans le plan \inserted{affine} \inserted{défini \unsure{par}} $\operatorname{Tr}u=\deleted{\unsure{2}}1$, engendrée par le \unsure{sommet} $\deleted{\unsure{\rho_s,\rho_a}}\rho_s$. \deleted{En c\unsure{ar.}}

% loc: left, diagonal at 7–14
% slanted
\margin{les sommets sont \uline{les} \uline{unipotents} les arêtes sont les paires d'unip. \unsure{dits conjugués|dits équivalents} : $(\rho_s,\underset{-\sigma\rho_s\sigma}{\underbrace{\overset{\rho_{s'},\,\text{où }s'=\sigma(s)}{\overbrace{\sigma\rho_s\sigma^{-1}}}}})$. i.e. paires d'unip\deleted{\unsure{ot}}otents $(u,u')$ tels que $\operatorname{Tr}uu'=-1$}

% loc: 9–16
En réduisant mod $n$, je dis que ça se factorise par $\mathrm{SL}(2,\mathbb Z/n\mathbb Z)/{\pm1}$ (tautologique, tel que j'ai présenté les choses), et la représentation obtenue de $C_n$ est celle que j'ai décrit tantôt (mais qui doit garder un sens, sans supposer $n$ premier $\ldots$)
% archive sheet number: 63

